% !TEX root =  main.tex
\chapter{Coherence Phenomenon}
\hspace{\parindent}
Light sources play significant roles on various optical processes and systems. Lights from known sources like lasers are nearly monochromatic and are known to be relatively coherent, while light from the sun is incoherent. The generated light beams from light sources of various degrees of coherence exhibit behaviors that are very different from one another. 

Coherence refers to the amount of correlation in the optical field at separate times or separate points within a beam of light \cite{Voelzbook}. High coherence leads to stationary interference effects such as the “fringing” and “ringing” structures that appeared in the simulated diffraction and imaging results for monochromatic (coherent) light \cite{Voelzbook}.

To be able to discuss coherence, it is convenient to first discuss classify coherence into two: temporal and spatial coherence. 

\section{Interference}
\hspace{\parindent} 




\section{Spatial Coherence}



\section{Temporal Coherence}



\section{Coherent Light}






\section{Incoherent Light}







\section{Partially-coherent Light}